% Cover letter for the IEEE TNNLS submission of Paper 2. Seen by the editors only (not by reviewers), so it names the authors.
% Claims follow the abstract's binding wording (SAP A6.12 branch A, A4.2.3); numbers are those printed in the abstract.
\documentclass[11pt]{article}
\usepackage[a4paper,margin=2.3cm]{geometry}
\usepackage{times}
\usepackage[hidelinks]{hyperref}
\setlength{\parindent}{0pt}
\setlength{\parskip}{6pt}
\pagestyle{empty}
\begin{document}

\begin{flushright}
School of Computer Science and Engineering\\
Vellore Institute of Technology\\
Vellore 632014, Tamil Nadu, India\\[4pt]
\today
\end{flushright}

The Editor-in-Chief\\
IEEE Transactions on Neural Networks and Learning Systems

Dear Editor-in-Chief,

We submit the manuscript \emph{``Learned Reliability-Aware Fusion of Token Evidence Raised Image-Text Verification Accuracy Under
Gaussian Noise in a Unified Masked Diffusion Model: A Preregistered Confirmation in UniDisc''} for consideration as a \textbf{Full
Paper} in IEEE Transactions on Neural Networks and Learning Systems.

Unified masked diffusion models generate and understand text and images with one network, and with the image tokens masked and a
caption visible, the same network assigns a probability to each observed image code. The manuscript reports two separately
preregistered studies of image-text verification under image damage in one such pretrained model (UniDisc, 1.4 billion parameters)
on Common Objects in Context image-caption pairs.
\begin{itemize}\setlength{\itemsep}{2pt}
\item In the first study, a gate discards low-likelihood image tokens. In the preregistered test no improvement was detected under
JPEG compression at quality 10, Gaussian noise with standard deviation 0.1, and a 128-pixel splice; for each tested corruption the
90\% confidence interval of the net accuracy change lay inside $\pm 0.05$ (500 pairs each), and the paper states that this does not
show that the gate has no effect.
\item In the second study, on a fresh preregistered fold, a frozen learned reliability-aware fusion of token evidence raised
image-direction verification accuracy under Gaussian noise with standard deviation 0.1 from 0.733 to 0.802 (net 0.069, 95\%
confidence interval 0.029 to 0.109; exact one-sided McNemar test, $p = 0.0005$), including with length-matched negatives, on 450
images not used to design or fit the weighting, at that one severity.
\end{itemize}
Both studies were fixed in time-stamped preregistrations before their outcomes were seen, every number in the paper is traced to
committed result files, and the limits of scope (one model, one dataset, one confirmed severity) are stated in the paper. We believe
the work fits the journal's scope in learning systems and information fusion, and that its preregistered design and its separation
of confirmed from exploratory results will interest the readership.

The manuscript follows the journal's double-anonymous review policy: the review copy contains no author names, affiliations,
funding or acknowledgment of the institution, and the code and data are available to reviewers in anonymized form. The
supplementary material is provided as a separate file. We confirm that this work is original, has not been published elsewhere,
and is not under review at any other publication. The authors declare no competing interests. All authors have approved the
manuscript and agree with its submission.

Thank you for considering our manuscript.

Sincerely,

Subramaniyaswamy Vairavasundaram (corresponding author)\\
School of Computer Science and Engineering, Vellore Institute of Technology, Vellore 632014, Tamil Nadu, India\\
E-mail: \href{mailto:subramaniyaswamy.v@vit.ac.in}{subramaniyaswamy.v@vit.ac.in}\quad
ORCID: \href{https://orcid.org/0000-0001-5328-7672}{0000-0001-5328-7672}

on behalf of all authors:\\
Shlok Goenka (ORCID \href{https://orcid.org/0009-0002-8581-5657}{0009-0002-8581-5657}),
Subramaniyaswamy Vairavasundaram (ORCID \href{https://orcid.org/0000-0001-5328-7672}{0000-0001-5328-7672}), and
Balaji G N (ORCID \href{https://orcid.org/0000-0002-5346-0989}{0000-0002-5346-0989})

\end{document}
